%! AP Statistics — Unit 8, Lesson 8.3 — Lesson Plan
\documentclass[10pt]{article}
\usepackage{apstats-article}
\usepackage{apstats-boxes}

\newcommand{\UnitNumberName}{Unit 8: Inference for Categorical Data: Chi-Square \quad}
\newcommand{\LessonNumberName}{Lesson 8.3: Carrying Out a Chi-Square Test for Goodness of Fit}

\begin{document}
\begin{center}
    {\Large\bfseries \CourseName: \SchoolYear \\
    {\small\bfseries \UnitNumberName \LessonNumberName}}
\end{center}

\begin{tcolorbox}[colback=sky, colframe=navy, boxrule=0.9pt, arc=2mm,
    left=2mm, right=2mm, top=1.5mm, bottom=1.5mm]
    \textbf{Primary Objective:} Students will calculate the chi-square test statistic,
    find and interpret the p-value using a chi-square table, and state a complete
    conclusion for a goodness-of-fit test.
    (Big Ideas: VAR, DAT)
\end{tcolorbox}

\renewcommand{\arraystretch}{1.6}
\begin{skillbox}[Priority Ideas \& Skills]{goldbox}
\begin{minipage}[t]{0.48\linewidth}
    \textbf{AP Skills 3.E, 4.B, 4.E}
    \begin{itemize}
        \item \textbf{VAR-8.F (Skill 3.E):} Calculate the chi-square test statistic
              $\chi^2 = \sum (O - E)^2 / E$ with $df = k - 1$ (number of categories minus 1).
        \item \textbf{VAR-8.G (Skill 3.E):} Determine the p-value using a chi-square
              table or technology.
        \item \textbf{DAT-3.I (Skill 4.B):} Interpret the p-value as the probability,
              given $H_0$ is true, of obtaining a test statistic as or more extreme.
        \item \textbf{DAT-3.J (Skill 4.E):} Compare $p$ to $\alpha$; reject or fail to
              reject $H_0$; justify a claim about the population.
    \end{itemize}
\end{minipage}
\hfill
\begin{minipage}[t]{0.48\linewidth}
    \textbf{Key Understandings}
    \begin{itemize}
        \item Each $(O - E)^2 / E$ term measures how far one category's observed
              count is from its expected count (squared and scaled).
        \item A larger $\chi^2$ statistic provides more evidence against $H_0$.
        \item The p-value is \emph{always} right-tailed for chi-square tests.
        \item Failing to reject $H_0$ does \emph{not} prove $H_0$ is true; it means
              the data are consistent with the null model.
        \item Conclusions must include: direction (reject/fail to reject),
              comparison ($p$ vs.\ $\alpha$), and context.
    \end{itemize}
\end{minipage}
\end{skillbox}

\begin{skillbox}[{Vocabulary, Concepts, \& Theorems}]{greenbox}
\begin{minipage}[t]{0.48\linewidth}
    \begin{itemize}
        \item \textbf{Chi-square test statistic ($\chi^2$)} ---
              $\chi^2 = \sum \dfrac{(O - E)^2}{E}$; measures total deviation of
              observed counts from expected counts
        \item \textbf{Degrees of freedom (GOF)} --- $df = k - 1$, where $k$ is
              the number of categories
    \end{itemize}
\end{minipage}
\hfill
\begin{minipage}[t]{0.48\linewidth}
    \begin{itemize}
        \item \textbf{P-value (chi-square)} --- area to the right of the observed
              $\chi^2$ under the chi-square distribution with $df = k - 1$
        \item \textbf{P-value interpretation} --- the probability, given $H_0$ is
              true, of getting a $\chi^2$ statistic as or more extreme than observed
        \item \textbf{Conclusion template} --- compare $p$ to $\alpha$; reject or
              fail to reject $H_0$; answer the research question in context
    \end{itemize}
\end{minipage}
\end{skillbox}

\begin{skillbox}[Activate Prior Knowledge \& Spiral Review]{sky}
\begin{minipage}[t]{0.48\linewidth}
    \textbf{Prior Knowledge Check (3--4 min)}
    \begin{itemize}
        \item Lesson 8.2: stating $H_0$ and $H_a$ for a GOF test;
              calculating expected counts ($E = n \cdot p_0$);
              verifying the large-counts condition (all $E \ge 5$).
        \item Unit 6: p-value interpretation (``given $H_0$ is true\ldots'')
              and conclusion language (``reject / fail to reject'').
        \item Connect: ``We set up the test last time. Today we carry it out.''
    \end{itemize}
\end{minipage}
\hfill
\begin{minipage}[t]{0.48\linewidth}
    \textbf{Connections Forward}
    \begin{itemize}
        \item Chi-square test for homogeneity $\to$ Lesson 8.4
        \item Chi-square test for independence $\to$ Lessons 8.5--8.6
        \item Two-way tables and conditional distributions $\to$ Lesson 8.4
        \item AP Exam free-response: always show formula, table, contributions,
              and full conclusion template
    \end{itemize}
\end{minipage}
\end{skillbox}

\begin{skillbox}[Hook \& Lesson]{sky}
\begin{minipage}[t]{0.48\linewidth}
    \textbf{Hook (3--5 min)}\\
    Display the botanist's 3:2:2:1 color claim. Ask: ``We set up the test last
    class. Now what? How do we actually decide whether the data contradict the
    claim?''\\[4pt]
    Show the observed counts side by side with the expected counts. Ask:
    ``Which category looks most surprising? How do we combine the deviations
    in all four categories into a single number?''\\[4pt]
    Reveal: the chi-square test statistic does exactly that --- it aggregates
    all $(O - E)^2 / E$ terms into one number that can be compared to a table.
\end{minipage}
\hfill
\begin{minipage}[t]{0.48\linewidth}
    \textbf{Lesson Flow (15--18 min)}
    \begin{enumerate}
        \item Introduce the formula: $\chi^2 = \sum (O - E)^2 / E$.
              Explain why we square: deviations above and below $E$ both count;
              why we divide by $E$: a deviation of 5 in a category with $E = 10$
              is more surprising than $E = 100$.
        \item Walk through the botanist example. Fill in the table row by row;
              compute each $(O - E)^2 / E$ contribution; sum to $\chi^2 \approx 3.367$.
        \item Degrees of freedom: $df = k - 1 = 4 - 1 = 3$. Why $k - 1$?
              Once $k - 1$ categories are fixed, the last is determined (counts
              must sum to $n$).
        \item Chi-square table: locate $df = 3$ row; find where 3.367 falls.
              The p-value $\approx 0.338$. Remind: always right-tailed.
        \item Interpretation and conclusion templates.
    \end{enumerate}
\end{minipage}
\end{skillbox}

\begin{skillbox}[Explicit Instruction \& Active Monitoring]{sky}
\begin{minipage}[t]{0.48\linewidth}
    \textbf{Direct Instruction (8 min)}
    \begin{itemize}
        \item Project the botanist table; call on students to supply each
              $(O - E)^2 / E$ contribution before revealing.
        \item Emphasize the p-value template: ``The probability, given the null
              hypothesis (3:2:2:1 ratio) is true, of obtaining a chi-square
              statistic of 3.367 or larger is approximately 0.338.''
        \item Model the full conclusion: ``Because $p = 0.338 > \alpha = 0.05$,
              we fail to reject $H_0$. We do not have convincing evidence that
              the color distribution of these flowers differs from the
              3:2:2:1 ratio.''
        \item Stress: ``convincing evidence'' (not ``proof'').
    \end{itemize}
\end{minipage}
\hfill
\begin{minipage}[t]{0.48\linewidth}
    \textbf{Active Monitoring}
    \begin{itemize}
        \item Watch for students adding signed residuals $(O - E)$ rather than
              squaring --- redirect with ``Why must the statistic be non-negative?''
        \item Cold-call: ``Which category contributed the most to $\chi^2$?
              What does that tell us?'' (Yellow: observed 13, expected 20 ---
              biggest discrepancy.)
        \item Check that students use $df = k - 1$, not $k$.
        \item Listen for misstatements of the p-value: ``the probability
              $H_0$ is true.'' Correct immediately.
    \end{itemize}
\end{minipage}
\end{skillbox}

\begin{skillbox}[Group Work \& Differentiation]{redbox}
\begin{minipage}[t]{0.48\linewidth}
    \textbf{Partner/Group Activity (12--15 min)}\\
    MLB Birth Month Distribution: Is birth month uniformly distributed
    for major-league players? Students carry out the full GOF test from
    hypotheses through conclusion using $n = 180$, $k = 12$ months.

    \begin{enumerate}
        \item State $H_0$ and $H_a$.
        \item Calculate expected counts; verify large-counts condition.
        \item Compute all 12 chi-square contributions; sum to $\chi^2$.
        \item Identify $df$; use table to bound/find the p-value.
        \item Write a complete conclusion at $\alpha = 0.05$.
    \end{enumerate}
\end{minipage}
\hfill
\begin{minipage}[t]{0.48\linewidth}
    \textbf{Differentiation}
    \begin{itemize}
        \item \textit{Support (Tier R)}: Provide the expected count (15) and the
              formula; students fill in the contribution column and sum.
        \item \textit{On-grade (Tier A)}: Full problem from scratch.
        \item \textit{Extension (Tier E)}: After finishing, identify which two
              months had the highest contributions. Hypothesize why (relative-age
              effect in youth sports); discuss whether the test result changes
              if December and January are combined into one category.
        \item \textit{ELL}: Bilingual vocabulary card ---
              ``test statistic / estadístico de prueba,''
              ``p-value / valor p,''
              ``fail to reject / no rechazar.''
    \end{itemize}
\end{minipage}
\end{skillbox}

\begin{skillbox}[Individual Work \& Assessment]{redbox}
\begin{minipage}[t]{0.48\linewidth}
    \textbf{Exit Ticket (5 min)}\\
    Candy color distribution: bag of 160 candies claims 25\% each of 4 colors.
    Observed counts given; chi-square contributions given; $\chi^2 = 2.0$,
    $df = 3$, $p > 0.10$.
    \begin{enumerate}
        \item Interpret the p-value in context.
        \item State a complete conclusion at $\alpha = 0.05$.
    \end{enumerate}
\end{minipage}
\hfill
\begin{minipage}[t]{0.48\linewidth}
    \textbf{AP-Style Check}\\
    \textit{A chi-square GOF test yields $\chi^2 = 7.81$ with $df = 3$.
    Which of the following is the best interpretation of the p-value
    if $p \approx 0.05$?}
    \begin{enumerate}[label=(\Alph*)]
        \item There is a 5\% probability the null hypothesis is true.
        \item If $H_0$ is true, about 5\% of samples give $\chi^2 \ge 7.81$.
        \item There is a 95\% probability the color distribution matches $H_0$.
        \item The data prove the distribution differs from $H_0$.
    \end{enumerate}
    Answer: (B)
\end{minipage}
\end{skillbox}

\begin{skillbox}[Reinforcement \& Extension]{goldbox}
\begin{minipage}[t]{0.48\linewidth}
    \textbf{Homework / Practice}
    \begin{itemize}
        \item Problem 1 (genetics 9:3:3:1 ratio, $n = 160$): calculate expected
              counts, compute $\chi^2$, identify $df$, state conclusion
              ($p \approx 0.81$, fail to reject).
        \item Problem 2 (streaming services, $n = 200$, 5 equal categories):
              state hypotheses, compute $\chi^2 \approx 53.65$, interpret
              p-value $< 0.001$, state conclusion.
        \item Problem 3 (conceptual): correct a student's error ---
              ``failing to reject $H_0$ proves the distribution matches.''
    \end{itemize}
\end{minipage}
\hfill
\begin{minipage}[t]{0.48\linewidth}
    \textbf{Extension / Preview}
    \begin{itemize}
        \item \textit{Extension}: Find a published data set (e.g., M\&M color
              proportions from Mars' website). Set up and carry out a full GOF
              test. Write a one-paragraph report interpreting the result.
        \item \textit{Preview}: Lesson 8.4 extends the chi-square procedure to
              compare distributions across two or more groups --- the
              chi-square test for homogeneity. Key question to bring to class:
              ``What changes when we have a two-way table instead of one row
              of counts?''
    \end{itemize}
\end{minipage}
\end{skillbox}

\end{document}
